\subsection{The ``up-to'' techniques}
\label{ss:upto}
The ``up-to'' techniques allow us to reduce
the size of a relation $\R$ to exhibit for proving bisimilarities.

\subsubsection{Bisimulation up-to context and expansion}

Our main up-to technique will be
{\em up-to context and expansion} \cite{San93c},
which admits the use of contexts and of behavioural equivalences such as 
expansion to achieve the
closure of a relation in the bisimulation game.

For our purposes, a restricted form of 'up-to context' will be enough,
in which only contexts of the form $\res \tilc (R   |  \contexthole)$ are
employed.


\begin{definition}[bisimulation up-to context and up-to $\contr$]
  \label{d:uptp_sce}
  A symmetric relation $  \R   \subseteq \pr \times \pr $ is a 
  {\em bisimulation up-to context and up-to $\contr$ }
  if $P \RR Q$ and 
  $P \arr\mu P''$
  then there are processes  $P'$ and $Q'$ 
  and a static context $\qct$
  s.t. $ P'' \contr \ct {P'}$, 
  $Q \Arcap\mu  \contr \ct{Q'}$ and 
  $P'  \RR  Q'$.
\end{definition}


The soundness of the above technique relies on the following lemma.
\begin{lemma}
\label{l:uptocon-expa}
Suppose $ \R $ is a bisimulation up-to context and up-to $\contr$,
 $(P, Q) \in  \R $ holds and  $\qct$ is a 
 static context. 
If  $\ct{P}\arr{\mu_1} \cdots\arr{\mu_n}  P_1$, $n\geq 0$,
 then
there are a static context $\qctp$ and processes $P'$ and $Q'$
s.t.\  $P_1 \contr \ctp{P'}$, $\ct Q \Arcap{\mu_1} \cdots
\Arcap{\mu_n} \contr \ctp{Q'}$ and $P'  \RR  Q'$. 
\end{lemma} 

\begin{theorem}
\label{t:up-to-con-expa}
If $ \R  $  is a 
bisimulation up-to context and up-to $\contr$, then $ \R  \subseteq \wbPI$.
\end{theorem} 
 
\begin{proof}
By showing that  
the relation $\S$ containing 
all pairs $(P,Q)$ such that 
\[ \begin{array}{l}
%  \R \defi \{ (P,Q) \st & 
 \mbox{ for some static context $\qct$ and
processes $P', Q'$} \\
%&
 \mbox{ it holds that } P \contr \ct{P'}, Q \contr \ct{Q'} \mbox{ and }
P'  \RR  Q' %&  \}
 \end{array}
 \]
is a bisimulation. 
See \cite{San93c} for details (adapting the proof to \Intp is
straightforward).
\end{proof}

\subsubsection{Bisimulation up-to bisimilarity}
\label{ss:bisimulation_upto_wb}

We will also employ {\em bisimulation up-to $\wbPI$}~\cite{Mil89}, 
a widely-used technique,
whereby bisimilarity itself is employed to achieve 
the closure of the candidate relation during the bisimulation game.


\begin{definition}
  [bisimulation up-to $\wbPI$]
  A symmetric relation $\R$ on $\pr  \times \pr$
  is a {\em bisimulation up-to $\wbPI$ } if $P  \RR  Q$ and $P
  \ArrPI\mu P'$ imply that  there exists $Q'$ s.t.
  $Q \Arcap\mu Q'$ and 
  $P' \wbPI  \RR \wbPI Q'$. 
\end{definition}
 

\begin{theorem}
\label{t:up-to}
  Suppose $\R$ is a {\em bisimulation up-to $\wbPI$}.
  Then $\R \subseteq \wbPI$.
\end{theorem}


\subsubsection{Bisimulation up-to context and up-to $(\contr,
\wbPI)$}

We will also employ
a variant of bisimulation up-to context and expansion, called 
{\em bisimulation up-to context and up-to $(\contr,\wbPI)$}.
In Definition~\ref{d:uptp_sce}, the occurrence of $\contr$  in 
$Q \Arcap\mu  \contr  \ct{Q'}$  could be weakened to $\wbPI$, and the
soundness of the technique, namely Theorem \ref{t:up-to-con-expa},
would still hold. However, in general, the same weakening
cannot be made on the requirement $P'' \contr \ct {P'}$.
We show that however, the weakening is possible when 
the action $\mu$ is visible and the context cannot interact with the
processes in its holes.

This proof technique, 
(which, as far as we know, does not appear in the literature) 
will be essential in some of our key results 
about the correspondence with game semantics.

\begin{definition}
  \label{d:interact}
  A static context $\qct$  \emph{cannot interact} with a process $P$
  if no name appears free both in the context and in
  the process and with opposite polarities (that is, either in input
  position in the context and output position in the process, or the
  converse).  
\end{definition} 

If  $\qct$ {cannot interact} with $P$,
whenever $\ct P \longrightarrow Q$ then  the
interaction is either internal to $\qct$ or internal to $P$. Moreover,
the property is invariant for reduction; i.e., 
if $Q = \ctp{P'}$ where $\qctp$ is the derivative of $\qct$ and $P'$
the derivative of $P$, then also 
$\qctp$  cannot interact with  $P'$.
Notice the invariance is false in the full $\pi$-calculus. 


We recall that we use $\ell$ to range over visible actions. 
\begin{definition}[bisimulation up-to context and up-to $(\contr,
 \wbPI)$]
\label{d:uptocontext_contr_wb}
A symmetric relation $  \R   \subseteq \pr \times \pr $ is a 
{\em bisimulation up-to context and up-to $(\contr,\wbPI)$}
if $P \RR Q$ imply:
  \begin{enumerate}
  \item if $P \arr\ell P''$
  then there are processes $P'$ and $Q'$ 
  and a static context $\qct$  
  s.t.\  $P'' \wbPI \ct {P'}$, $Q \Arr\ell  \wbPI  \ct{Q'}$ 
  and $P' \RR Q'$, 
  and, moreover, $\qct$ cannot interact with $P$ or $Q$; 
  %%
  \item if $P \longrightarrow P''$
  then there are processes $P'$ and $Q'$
  and a static context $\qct$
  s.t.\  $ P'' \contr \ct {P'}$, 
  $Q \Longrightarrow \contr  \ct{Q'}$ 
  and $P' \RR Q'$.
  \end{enumerate}
\end{definition}


\begin{lemma}
  \label{l:upto_new}
  Suppose $ \R  $  is a 
  bisimulation up-to context and up-to $(\contr,\wbPI)$ 
  and $\qct$ is a static context 
  that does not interact with processes $P$ and $Q$.  
  \begin{enumerate}
  \item  
  If  $\ct{P}\Longrightarrow   P_1$, 
  then there are processes $P'$ and $Q'$ 
  and a static context $\qctp$
  s.t.\  
  $P_1 \contr \ctp{P'}$, 
  $\ct Q \Longrightarrow \contr \ctp{Q'}$ 
  and $P'  \R  Q'$. 
  
  \item 
  If $\ct{P}\arr\ell   P_1$, 
  then there are processes $P'$ and $Q'$ and 
  a static context $\qctp$ and 
  s.t.\  
  $P_1 \wbPI \ctp{P'}$, 
  $\ct Q \ArrPI\ell \wbPI \ctp{Q'}$ 
  and $P'  \R  Q'$. 
  \end{enumerate} 
\end{lemma}

\noindent
The proof of  Lemma~\ref{l:upto_new} is along the lines of similar results in the
literature combining up-to expansion and up-to bisimilarity, see e.g.~\cite{SaWabook}. 

\reviewComment{there is no proof. Is that a classical lemma? Is it
because the proof is straightforward?}
\reviewCommentDS{Answer: while strictly speaking the result does not appear in the
  literature, the proof is straightforward adapation of  similar techniques. We have added
a sentence with a reference  (To be honest: the result is not entirely trivial, and
``normally'' one should have indeed added a proof. In other words, the referee  is right
is complaining. However it is certainly not the most important proof that is missing from
the main text! This paper should have really been 2 or even 3 ordinary papers)}

\begin{theorem}
\label{t:up-to-con-expa-nointer} 
  If $\R $ is a bisimulation 
  up-to context and up-to $(\contr,\wbPI)$, 
  then ${ \R } \subseteq {\wbPI}$.
\end{theorem}

\begin{proof}
  By showing that 
  the relation $\R$ containing 
  all pairs $(P,Q)$ such that 
  \[
   \begin{array}{l}
  %  \R \defi \{ (P,Q) \st & 
   \mbox{ for some static context $\qct$ and
  processes $P', Q'$} \\
  %&
   \mbox{ where $\qct$ does not interact with $P',Q'$ }  \\
  %& 
  \mbox{ it holds that } P \wbPI \ct{P'}, Q \wbPI \ct{Q'} \mbox{ and }
  P'  \RR  Q' %&  \}
   \end{array}
   \]
  is a bisimulation. 
  In a challenge transition $P \arr\mu P_1$ one distinguishes the case
  when $\mu$ is a silent or visible action. 
  In both cases, one exploits Lemma~\ref{l:upto_new}. 
\end{proof}
